\setstretch{1.5}
\setlength{\parskip}{0.5\baselineskip}

This chapter builds the background necessary to follow the rest of the thesis.
A reader unfamiliar with Arabic script, OCR technology, large language models,
or prompt engineering should be able to acquire the required context here.
Section~\ref{sec:arabic} covers the Arabic language and script.
Section~\ref{sec:ocr} introduces optical character recognition and its
Arabic-specific challenges.
Section~\ref{sec:llms} surveys large language models from early neural approaches
through the modern transformer era.
Section~\ref{sec:prompting} covers prompt engineering methods relevant to this
study.
Section~\ref{sec:knowledge_sources} describes the external knowledge sources
used.
Section~\ref{sec:related} reviews related work on post-OCR correction and
positions this thesis within that literature.

% ============================================================
\section{The Arabic Language and Script}
\label{sec:arabic}
% ============================================================

\subsection{Linguistic Overview}
\label{subsec:arabic_linguistic}

Arabic is a Semitic language and the liturgical language of Islam, spoken by
over 300 million people as a first language across the Arab world and
understood by a much larger community as a second or religious language
\cite{holes2004modern}.
Its most distinctive grammatical feature is \emph{root-and-pattern morphology}:
the vast majority of Arabic words are derived from a three-consonant (trilateral)
root, and different word patterns are created by interleaving the root consonants
with fixed vowel-and-consonant patterns.
For example, the root \textit{k-t-b} (relating to writing) yields:
\textit{kataba} (he wrote), \textit{kitaab} (book), \textit{kaatib} (writer),
\textit{maktaba} (library), \textit{maktub} (written/letter), and dozens of
other morphologically related forms.

This morphological richness creates a large and sparse vocabulary for text
processing tools.
Where English uses analytic constructions (``he will have been writing''),
Arabic encodes tense, aspect, voice, gender, number, and definiteness
through bound morphology, generating many distinct surface forms for each
lemma.
A vocabulary estimated at several hundred thousand distinct word forms makes
statistical NLP methods data-hungry and challenges neural models trained on
smaller corpora.

Modern Standard Arabic (MSA) is the standardised written form used in formal
contexts across the Arab world: newspapers, government documents, academic
publications, and formal broadcasts.
It descends from Classical Arabic (the language of the Quran) with some
modernisation of vocabulary.
Alongside MSA, a spectrum of spoken dialects exists (Egyptian, Gulf, Levantine,
Maghrebi, etc.) that differ substantially from MSA and from each other in
phonology, morphology, and lexicon.
This study works with MSA throughout, as PATS-A01 documents are printed in MSA
and the QALB and OpenITI corpora also contain MSA text.

\subsection{The Arabic Script}
\label{subsec:arabic_script}

Arabic is written with a 28-letter abjad (consonantal alphabet) that runs
from right to left.
Several properties of the script are critical for understanding OCR failure modes.

\textbf{Positional forms.}
Most letters have four shapes depending on their position within a word:
isolated, initial, medial, and final.
The letter \textit{ain} (\textarabic{�}), for instance, appears as
\textarabic{�} (isolated), \textarabic{��} (initial), \textarabic{���} (medial),
and \textarabic{��} (final).
Since letters are connected into words by ligatures (joined strokes), an OCR
engine must recognise not just a character in isolation but a glyph in context,
making segmentation errors compound recognition errors.

\textbf{Diacritics (tashkeel / harakat).}
Short vowels in Arabic are written as small marks placed above or below
consonant letters.
The primary marks are: fatha (short /a/, above), kasra (short /i/, below),
damma (short /u/, above), sukun (consonant with no vowel, above), and shadda
(gemination, above).
Tanwin (nunation) suffixes -- tanwin al-fath, tanwin al-kasr, tanwin al-damm --
mark indefinite nouns and consist of doubled fatha, kasra, or damma marks.
In most modern Arabic text, including newspapers, diacritics are \emph{omitted}
from printed text and readers infer the correct vowelisation from context.
When present, diacritics are critical for disambiguation: the same consonant
sequence can represent multiple words depending on vowelisation.

\textbf{Hamza.}
The glottal stop (hamza) has a complex orthography.
In some positions it appears alone (\textarabic{�}); in others it appears on a
carrier letter: alef (\textarabic{�} above, \textarabic{�} below), waw (\textarabic{�}),
or ya without dots (\textarabic{�}).
The choice of carrier follows rules based on the surrounding vowels, and even
educated native speakers sometimes apply these rules inconsistently.
Hamza rules are a major source of spelling errors in human writing and a
significant challenge for OCR.

\textbf{Taa marbuta.}
The feminine ending \textarabic{�} (taa marbuta) is pronounced /at/ in context
and /a/ in pause.
It is written with the same base shape as the letter ha (\textarabic{�}) but
with two dots above.
The visual similarity between \textarabic{�} and \textarabic{�} is a direct
source of OCR errors whenever the ink quality, resolution, or font rendering
is imperfect.

\textbf{Alef maqsura.}
The letter \textarabic{�} (alef maqsura, literally ``abbreviated alef'') appears
at the end of certain words to write the long vowel /aa/.
Its shape is identical to ya (\textarabic{�}) but without the two dots below.
OCR engines regularly confuse \textarabic{�} and \textarabic{�}, and the
distinction is orthographically important because the two letters signal
different word identities.

\textbf{Dot-group characters.}
Six groups of letters share the same base skeletal shape and differ only in
the count or position of diacritical dots above or below:
\textarabic{�}~/~\textarabic{�}~/~\textarabic{�}~(1/2/3 dots above),
\textarabic{�}~/~\textarabic{�}~/~\textarabic{�}~(1 dot below / none / 1 dot above),
\textarabic{�}~/~\textarabic{�}~(1 dot above / 2 dots above),
\textarabic{�}~/~\textarabic{�}~(1 dot above / 2 dots below), and others.
Dot placement errors -- caused by ink smearing, low resolution, or font
rendering artefacts -- are the most frequent single-character OCR confusion
in Arabic typewritten text.

\subsection{Arabic Text in Computing}
\label{subsec:arabic_computing}

Arabic text in digital systems is encoded in Unicode.
The primary Arabic Unicode block spans U+0600 to U+06FF and covers the standard
28 letters, their positional variants, diacritical marks, punctuation, and
numerals.
The diacritical marks (tashkeel) occupy code points U+064B through U+065F for
the standard marks, with additional marks in the ranges U+0610--U+061A and
U+0670.
These code points can be removed programmatically to strip diacritics without
affecting the consonantal skeleton.

Right-to-left rendering requires bidirectional text support.
The Unicode Bidirectional Algorithm (UBA) governs how mixed LTR/RTL content is
laid out.
Documents containing both Arabic and numerals (which are LTR) or Arabic and
English words exhibit complex rendering behaviour that can cause misalignment
between stored byte order and visual display order.
OCR engines and NLP tools must handle this correctly to avoid misinterpreting
character sequences.

Legacy encodings such as Windows Code Page 1256 (CP-1256), ISO-8859-6, and the
ASMO-708 standard predate Unicode and are still found in digitised Arabic
document collections.
The PATS-A01 ground-truth files, for example, are stored in CP-1256.
Correct decoding requires explicit encoding specification; reading an Arabic
CP-1256 file as UTF-8 produces garbled output.
All text in this study is converted to UTF-8 on load.

\subsection{Challenges for Arabic NLP}
\label{subsec:arabic_nlp_challenges}

\textbf{Tokenisation.}
Arabic word boundaries are clear (words are separated by spaces), but
morpheme boundaries within words are not.
Arabic clitics -- conjunctions, prepositions, definite articles, pronominal
suffixes -- are written attached to the host word without any separator.
Tokenising the sentence into morphological units requires a separate
segmentation step beyond whitespace splitting.

\textbf{Vocabulary sparsity.}
The large number of inflectional and derivational forms per lemma means that
any fixed-vocabulary system encounters out-of-vocabulary tokens frequently
unless the vocabulary is very large.
Subword tokenisation approaches (byte-pair encoding, WordPiece) partially
address this but fragment morphologically complex words into unintuitive
subword units.

\textbf{Dialectal variation.}
Egyptian, Gulf, Levantine, and Maghrebi dialects differ from MSA and from each
other in phonology, morphology, and lexicon.
A model trained on MSA may perform poorly on dialectal text.
This study focuses on MSA text throughout.

\textbf{Code-switching.}
Modern Arabic digital text frequently mixes Arabic script with embedded
English words, numbers, emoticons, and URLs.
Handling such mixed-script content robustly requires robust character-level
processing.

% ============================================================
\section{Optical Character Recognition}
\label{sec:ocr}
% ============================================================

\subsection{What Is OCR?}
\label{subsec:ocr_what}

Optical Character Recognition (OCR) is the process of converting images of
text into machine-readable character sequences.
The classical OCR pipeline consists of four stages:

\begin{enumerate}
  \item \textbf{Preprocessing.}
        The scanned or photographed document image is prepared for analysis:
        binarisation (converting to black-and-white), deskewing (correcting
        rotational distortion), denoising (removing salt-and-pepper and
        background noise), and normalisation of brightness and contrast.

  \item \textbf{Layout analysis.}
        The preprocessed image is segmented into regions: page borders, columns,
        paragraphs, text lines, and (in character-segmenting systems) individual
        character images.
        For Arabic, line detection must handle the right-to-left reading order
        and the slanted or curved baselines that appear in handwritten text.

  \item \textbf{Character recognition.}
        The character recognition module maps each segmented glyph (or a sequence
        of glyphs, in modern end-to-end systems) to a Unicode character.
        This is the core of OCR and the component that has changed most
        dramatically with deep learning.

  \item \textbf{Post-processing.}
        A language model or lexicon is consulted to resolve ambiguous recognitions
        and correct plausible errors.
        In classical systems this was a rule-based or n-gram language model step.
        In modern systems it is sometimes omitted because the recognition model
        is end-to-end and incorporates linguistic knowledge implicitly.
\end{enumerate}

\subsection{OCR Architectures}
\label{subsec:ocr_arch}

The OCR field has passed through three distinct architectural eras:

\textbf{Template matching (pre-2000).}
Early systems compared each segmented character image against a library of
stored templates and selected the nearest match by pixel distance or correlation
\cite{mori1992historical}.
Template matching is brittle with respect to font variation, noise, and
degradation, and requires re-training the template library for each new font.

\textbf{Hidden Markov Models (1990s--2010s).}
HMM-based systems model the sequence of observations (image features along a
text line) as a probabilistic generative process.
They can handle cursive scripts without explicit character segmentation, which
made them the dominant approach for Arabic OCR through the 2000s.
A character language model is typically combined with the HMM in a noisy channel
framework to improve recognition.

\textbf{Deep learning era (2014--present).}
Convolutional neural networks replaced hand-crafted feature extraction.
The combination of CNN feature extraction, Bidirectional LSTM sequence modelling,
and Connectionist Temporal Classification (CTC) decoding \cite{graves2006connectionist}
became the dominant paradigm for text-line recognition.
Tesseract v4, the open-source OCR engine on which Qaari is based, adopted this
CNN-LSTM-CTC architecture \cite{smith2007overview}.

\textbf{Transformer-based OCR (2021--present).}
Vision-language models such as TrOCR \cite{li2021trocr} and Donut
\cite{kim2022donut} treat OCR as a sequence-to-sequence generation task: a
vision encoder encodes the document image, and a text decoder generates the
recognised text token by token.
These models are pre-trained on large paired image-text datasets and achieve
state-of-the-art performance on printed and even handwritten text.
However, they require substantial compute and training data, which limits their
availability for low-resource languages and scripts.

\subsection{Arabic-Specific OCR Challenges}
\label{subsec:arabic_ocr_challenges}

The properties of the Arabic script described in Section~\ref{sec:arabic}
translate directly into OCR failure modes.

\textbf{Segmentation errors.}
Because Arabic letters are connected within a word, a segmentation error
(incorrectly splitting a ligature or incorrectly merging two letters) cascades
into recognition errors for all characters involved.
English OCR, which deals mostly with separated or lightly connected Latin
characters, does not face this problem to the same degree.

\textbf{Dot placement ambiguity.}
The fine dots above and below base letter shapes carry discriminative information
but are small relative to the letter body.
Low image resolution, ink bleed, or scan artefacts can render dots illegible,
and the OCR engine must then guess between all members of the relevant dot group.
This is the most frequent source of character-level errors in PATS-A01.

\textbf{Ligature handling.}
Certain letter combinations in Arabic merge into single ligature glyphs with
a shape that may not visually resemble the component letters in isolation.
Common examples include lam-alef (\textarabic{��}), which is always written as
a single glyph.
OCR engines trained on character-level templates may not include all ligature
shapes, leading to substitution or insertion errors.

\textbf{Diacritic sensitivity.}
OCR engines vary in whether they attempt to recognise diacritics.
Many ignore them entirely, which is appropriate for most modern Arabic text
but problematic for classical or religious text where diacritics are always
present.
Others attempt recognition but produce errors when diacritics are present in
unusual positions or rendered in non-standard fonts.

\textbf{Baseline variability.}
In handwritten text, words do not lie on a flat baseline.
Variable baseline trajectories challenge text-line segmentation and make
aligned feature extraction more difficult.

\subsection{Qaari: The OCR Engine Used in This Study}
\label{subsec:qaari}

Qaari is an open-source Arabic OCR engine based on Tesseract v4 with Arabic
LSTM models trained specifically for Arabic printed text.
It represents the class of freely available, community-maintained OCR engines
that are accessible to researchers and practitioners who cannot afford commercial
OCR solutions.

In the experiments of this thesis, Qaari performs well on clean PATS-A01 Akhbar
text: the character error rate on normal samples is 1.37\% (Section~\ref{ch_results}).
However, Qaari exhibits one notable failure mode: a \emph{repetition artifact}.
When processing certain document layouts, Qaari enters a loop in which it
repeatedly outputs the same short text span, producing a hypothesis that is
hundreds of times longer than the ground truth.
This artifact is rare (observed in 1 out of 200 samples, 0.5\%) but when it
occurs it produces an extremely high per-sample CER that dominates the aggregate.
The artifact is a real-world behaviour of the OCR engine and is included in all
evaluations to preserve ecological validity.

\subsection{Evaluation Metrics for OCR}
\label{subsec:ocr_metrics}

The two standard metrics for evaluating OCR and post-OCR correction are Character
Error Rate (CER) and Word Error Rate (WER), both defined in
Section~\ref{sec:eval} (Chapter~3).
Here we comment on the choice of CER as the primary metric and on the
diacritics stripping decision.

\textbf{Why CER over WER.}
Arabic morphology creates a large space of valid word forms; any character-level
error produces a word-level error.
WER is therefore always higher than CER for Arabic text and is more sensitive to
the tokenisation convention used.
CER captures the actual number of character edits needed, which is more directly
interpretable as a measure of recognition accuracy at the signal level.
CER is also the standard primary metric in post-OCR correction shared tasks
\cite{rigaud2019icdar}.

\textbf{Diacritics stripping.}
Evaluating with diacritics would penalise a corrector for failing to reproduce or
adding diacritics that are stylistically optional.
Since neither Qaari nor the LLM corrector is designed to hallucinate or remove
diacritics that are part of the ground truth, evaluating without diacritics
provides a cleaner measure of lexical correction quality.
Both variants are computed and stored; the diacritic-stripped metric is reported
as primary throughout this thesis.

% ============================================================
\section{Large Language Models}
\label{sec:llms}
% ============================================================

\subsection{Why LLMs for Post-OCR Correction}
\label{subsec:why_llms}

Before introducing the technical foundations of large language models, it is
worth establishing \emph{why} they are the right tool for Arabic post-OCR
correction, and why this thesis adopts a prompt-engineering paradigm rather
than the traditional or deep-learning approaches reviewed in
Section~\ref{sec:related}.

Post-OCR correction is fundamentally a \emph{conditional text restoration}
task: given a noisy observation of a clean text (the OCR output), produce the
most likely original.  Historically, this task has been addressed by three
paradigm families, each with characteristic limitations for Arabic.

\textbf{Rule-based and statistical methods} (n-gram language models, lexicon
lookup, noisy-channel formulations) require hand-crafted rules or
frequency tables for each error type and language pair.
For Arabic, the combinatorial explosion of positional glyph variants,
dot-group confusions, and morphological inflections makes rule curation
infeasible at the scale required for robust coverage.
N-gram models capture local context (typically 3--5 tokens) but cannot
leverage sentential or discourse-level coherence, which is essential for
resolving valid-but-wrong substitutions --- the dominant error type in Arabic
OCR (71\% of errors in this study).

\textbf{Neural sequence-to-sequence models} (encoder--decoder
architectures such as those based on LSTMs or smaller transformers)
treat correction as a translation task from noisy to clean text.
They can learn complex error patterns from data but require large parallel
corpora of (OCR output, ground truth) pairs.
For Arabic OCR, such corpora are scarce: the only substantial Arabic
error correction dataset (QALB~\cite{zaghouani2014qalb}) targets \emph{grammatical}
errors in human-authored text, not \emph{visual} errors produced by OCR
engines.
Building an OCR-specific training corpus requires rendering text in multiple
fonts, running OCR, and aligning output with ground truth --- a process
that is engine-specific (a model trained on Tesseract errors does not
transfer to Qaari errors) and must be repeated for each font and engine
combination.
Furthermore, seq2seq models trained on one error distribution are brittle
when the error profile changes: a font change alters the confusion matrix
and can invalidate the learned correction mapping.

\textbf{Large language models} (LLMs) offer a fundamentally different
proposition.
A model pre-trained on hundreds of billions of tokens of multilingual text
acquires, as a byproduct of next-token prediction, an implicit probability
distribution over valid Arabic sentences.
This distribution functions as an extremely rich language model that captures
not just n-gram statistics but morphological constraints, syntactic structure,
semantic coherence, and world knowledge --- all without any task-specific
training.
Instruction tuning then teaches the model to follow explicit correction
instructions, converting this passive language competence into an active
error-correction capability.

The advantages of this paradigm for Arabic post-OCR correction are:

\begin{enumerate}
  \item \textbf{No parallel training data required.}  The model's correction
        capability derives from its pre-trained language prior, not from
        OCR-specific training pairs.  This eliminates the data bottleneck that
        limits seq2seq approaches for low-resource OCR correction scenarios.

  \item \textbf{Engine and font agnosticism.}  Because the model corrects
        based on its understanding of valid Arabic rather than on memorised
        error patterns, the same model and prompt can be applied to output from
        any OCR engine without retraining.  Knowledge augmentation (Phases~3
        and~4) can adapt the prompt to a specific engine's error profile, but
        the base capability transfers.

  \item \textbf{Rich contextual reasoning.}  Unlike n-gram models, LLMs
        process the full input sequence through multi-layer self-attention,
        enabling them to use long-range context to disambiguate locally
        ambiguous tokens.  This is particularly important for Arabic, where
        valid-but-wrong OCR errors can only be detected by semantic or
        syntactic analysis of the surrounding sentence.

  \item \textbf{Zero-shot applicability.}  An instruction-tuned LLM can
        perform correction immediately upon deployment, without any
        task-specific training phase.  This enables rapid experimentation with
        different prompt strategies (the core methodology of this thesis)
        without the computational cost of model training.

  \item \textbf{Interpretable prompt engineering.}  The knowledge injected
        into the model's prompt (confusion matrices, self-reflective insights,
        correction rules) is human-readable and auditable.  Unlike neural
        model weights, prompt content can be inspected, modified, and
        version-controlled, supporting transparent and reproducible research.
\end{enumerate}

The primary risk of the LLM paradigm --- and a central finding of this
thesis --- is \emph{over-correction}: the model's tendency to rewrite
tokens that are already correct, introducing new errors.
This risk is inherent to instruction-following models operating on text that
is mostly correct, and it is the reason that knowledge augmentation
(Phases~3--5) and prompt design (conservative v2 prompt) are necessary to
constrain the model's intervention behaviour.
Understanding and quantifying this over-correction tendency is a major
contribution of this work.

\subsection{From N-grams to Transformers}
\label{subsec:lm_history}

Statistical language modelling has a long history in NLP.
N-gram language models assign a probability to a sequence of words by decomposing
the joint probability into a product of conditional probabilities of each word
given the preceding $n-1$ words:
\begin{equation}\label{eq:ngram}
  P(w_1, w_2, \ldots, w_m) \approx \prod_{i=1}^{m} P(w_i \mid w_{i-n+1}, \ldots, w_{i-1})
\end{equation}
They were the dominant approach for language modelling and speech recognition
through the 1990s and 2000s.
Their limitations --- sparse counts for long contexts (the \emph{curse of
dimensionality}: the number of possible n-grams grows exponentially with $n$),
inability to generalise across syntactically similar contexts that happen to use
different vocabulary, and the practical ceiling of $n \leq 5$ due to memory
constraints --- motivated neural approaches.

Bengio et al.\ \cite{bengio2003neural} introduced the neural probabilistic language
model, which represents each word as a dense real-valued vector (an
\emph{embedding}) in a continuous space and uses a feedforward network to predict
the next word given a fixed-width context window.
This model demonstrated that neural approaches could outperform n-gram models by
exploiting the geometric structure of embedding space: words used in similar
contexts are mapped to nearby vectors, enabling the model to generalise from
observed contexts to unseen but structurally similar ones.

Recurrent Neural Networks (RNNs) extended this to variable-length contexts by
maintaining a hidden state vector that is updated at each time step, enabling the
model to condition its prediction on the entire preceding sequence.
Long Short-Term Memory (LSTM) networks \cite{hochreiter1997long} addressed the
vanishing gradient problem of vanilla RNNs with gated memory cells --- input,
forget, and output gates that control the flow of information into and out of a
persistent cell state --- enabling effective modelling of long-range
dependencies over hundreds of tokens.
LSTM-based language models and encoder-decoder sequence-to-sequence models
\cite{sutskever2014sequence} dominated NLP in the early 2010s, achieving
state-of-the-art results in machine translation, text summarisation, and
dialogue systems.

The transformer architecture \cite{vaswani2017attention}, introduced in 2017,
replaced recurrence entirely with self-attention mechanisms.
By computing attention over all positions simultaneously rather than sequentially,
transformers eliminated the sequential bottleneck of RNNs and enabled
embarrassingly parallel training on modern GPU hardware.
The result was a dramatic improvement in training efficiency: transformers train
orders of magnitude faster than LSTMs on the same data and, critically, exhibit
favourable \emph{scaling laws} --- their performance improves predictably as a
power law of model size, dataset size, and compute
\cite{kaplan2020scaling,hoffmann2022training}.
This scalability property is the foundation of the modern LLM paradigm: given
sufficient data and compute, a transformer-based model can be scaled to sizes
where qualitatively new capabilities emerge --- instruction following, few-shot
learning, chain-of-thought reasoning --- that are not present at smaller scales
\cite{wei2022emergent}.
All modern large language models are transformer-based.

\subsection{The Transformer Architecture}
\label{subsec:transformer}

The key innovation of the transformer is the \emph{self-attention mechanism},
which allows every token in a sequence to directly attend to every other token,
regardless of their distance.
Given a sequence of token embeddings represented as a matrix $X \in \mathbb{R}^{n \times d}$,
the scaled dot-product attention is computed as:
\begin{equation}\label{eq:attention}
  \mathrm{Attention}(Q,K,V) = \mathrm{softmax}\!\left(\frac{QK^\top}{\sqrt{d_k}}\right) V
\end{equation}
where $Q = XW^Q$, $K = XW^K$, $V = XW^V$ are learned linear projections of the
input into query, key, and value spaces, and $d_k$ is the key dimension used for
scaling.
The softmax produces an attention weight distribution over all positions, and the
output is a weighted sum of value vectors.

\textbf{Multi-head attention} runs $h$ independent attention operations in parallel,
each with its own projection matrices, and concatenates their outputs:
\begin{equation}\label{eq:multihead}
  \mathrm{MultiHead}(Q,K,V) = \mathrm{Concat}(\mathrm{head}_1, \ldots, \mathrm{head}_h) W^O
\end{equation}
This allows the model to attend to information from different representational
subspaces simultaneously.

\textbf{Positional encoding} is added to the input embeddings to provide
information about token order, since the attention operation is otherwise
position-invariant.
Original transformers used sinusoidal positional encodings; modern LLMs often
use rotary positional embeddings (RoPE) or similar learned variants.

A transformer block consists of a multi-head self-attention sublayer followed
by a position-wise feedforward network, with residual connections and layer
normalisation around each sublayer.
Modern LLMs are decoder-only transformers: they use masked (causal) self-attention
so that each token can only attend to tokens at or before its position, enabling
autoregressive generation.

\subsection{Pre-Training and Instruction Tuning}
\label{subsec:pretraining}

Modern LLMs are trained in two or three stages, each adding a qualitatively
different capability.

\textbf{Stage 1: Pre-training (next-token prediction).}
The model is trained to predict the next token given all preceding tokens, on a
massive corpus of text scraped from the internet, books, code, scientific
papers, and other sources --- typically trillions of tokens spanning dozens
of languages.
The training objective is the standard autoregressive language modelling loss:
\begin{equation}\label{eq:pretraining}
  \mathcal{L} = -\sum_{t=1}^{T} \log P_\theta(x_t \mid x_1, \ldots, x_{t-1})
\end{equation}
where $\theta$ denotes the model parameters and $x_1, \ldots, x_T$ is a
tokenised training sequence.
Despite its simplicity, this objective is remarkably powerful: by learning to
predict the next token across a sufficiently large and diverse corpus, the model
acquires broad language competence, world knowledge, logical reasoning patterns,
and the ability to perform tasks it was never explicitly trained for --- purely
from the statistical structure of human language.
GPT-3 \cite{brown2020language} demonstrated that scaling pre-training to 175
billion parameters enables dramatic improvements in few-shot performance across
NLP tasks, establishing the empirical basis for the \emph{scaling hypothesis}:
larger models trained on more data exhibit qualitatively new capabilities.

For Arabic post-OCR correction, the critical consequence of pre-training is
that the model learns an implicit distribution over valid Arabic text.
When presented with a noisy OCR sequence, the model's next-token predictions
are biased toward tokens that form valid Arabic words and coherent sentences,
effectively functioning as a powerful contextual spell-checker --- but one that
operates at the semantic and syntactic level, not just the lexical level.

\textbf{Stage 2: Instruction tuning (supervised fine-tuning, SFT).}
A pre-trained model produces likely continuations of text; it does not
reliably follow explicit instructions.
Instruction tuning \cite{ouyang2022training} fine-tunes the model on a curated
set of (instruction, response) pairs, teaching it to interpret natural-language
instructions as tasks and produce helpful, structured outputs.
The fine-tuning dataset covers diverse tasks --- question answering, summarisation,
translation, code generation, text editing --- formatted as multi-turn dialogues
with system prompts, user messages, and assistant responses.
After instruction tuning, the model responds predictably to prompts of the form
``correct this OCR text and return only the corrected version,'' making it
directly applicable to the post-OCR correction task without further training.

\textbf{Stage 3: Alignment (RLHF / DPO).}
Reinforcement Learning from Human Feedback (RLHF)~\cite{ouyang2022training}
further aligns the model's outputs with human preferences.
A reward model is trained on human comparisons of model outputs (``output A is
better than output B for this prompt''), and the language model is then
optimised against this reward model using policy gradient methods (typically
Proximal Policy Optimisation, PPO~\cite{schulman2017proximal}).
More recent approaches such as Direct Preference Optimisation
(DPO)~\cite{rafailov2023direct} bypass the reward model entirely and
optimise the policy directly on preference pairs.
The result is a model that not only follows instructions but does so in a way
that is helpful, harmless, and honest --- reducing the frequency of
hallucinations, harmful outputs, and unnecessary rewrites.

For post-OCR correction, alignment is particularly important because it
reduces the model's tendency to ``hallucinate'' --- to generate plausible but
incorrect text that replaces what the OCR actually produced.
An aligned model is more likely to return the input unchanged when it is
already correct, which is precisely the conservative behaviour needed to
avoid over-correction on low-error input.

\subsection{Qwen3-4B-Instruct-2507}
\label{subsec:qwen3}

Qwen3-4B-Instruct-2507 is the primary LLM used in this study.
It belongs to the Qwen3 family released by Alibaba's Tongyi team \cite{qwen3report}.
Key characteristics relevant to this study are as follows.

\textbf{Scale and architecture.}
Qwen3-4B has approximately 4 billion parameters and is a decoder-only transformer.
At this scale it fits comfortably in the 16~GB VRAM of a Kaggle T4 GPU in
\texttt{float16} precision, making it a practical choice for cloud GPU inference
without quantization.

\textbf{Multilingual capability and Arabic support.}
Qwen3 was pre-trained on a large multilingual corpus covering over 100 languages,
with substantial Arabic representation.
This enables the model to read, generate, and reason in Arabic without any
Arabic-specific fine-tuning.

\textbf{Instruction following.}
The \texttt{-Instruct} variant is trained with supervised fine-tuning and
preference alignment to follow instructions reliably.
This is the essential property for post-OCR correction, where the model must
interpret a system prompt and produce a structured text output.

\textbf{Thinking mode.}
Qwen3 models include an optional chain-of-thought reasoning mode controlled by
the \texttt{enable\_thinking} parameter in \texttt{apply\_chat\_template()}.
When enabled, the model generates an extended internal reasoning trace enclosed
in \texttt{<think>...</think>} tokens before producing its final answer.
This study disables thinking mode (\texttt{enable\_thinking=False}) for two
reasons: first, post-OCR correction is a direct transformation task that does
not benefit from extended chain-of-thought reasoning; second, keeping thinking
mode enabled adds latency and token cost, and the thinking trace contents
sometimes contaminate the extracted correction.

\subsection{Arabic LLMs and Arabic NLP with LLMs}
\label{subsec:arabic_llms}

The Arabic NLP community has produced several dedicated pre-trained models.

\textbf{AraBERT} \cite{antoun2020arabert} is an encoder-only BERT model
pre-trained on a large corpus of Arabic news and Wikipedia text.
It achieves strong results on Arabic classification and understanding tasks
and established that pre-training specifically on Arabic outperforms
multilingual BERT on Arabic benchmarks at comparable model size.

\textbf{CAMeLBERT} \cite{inoue2021camelbert} extends this line of work
with models pre-trained on MSA, dialectal Arabic, and mixed-domain corpora,
providing tools tailored to different Arabic registers.

\textbf{AraGPT2} \cite{antoun2021aragpt2} and \textbf{AraT5}
\cite{nagoudi2021arat5} bring generative and text-to-text capabilities to Arabic,
enabling Arabic language generation and conditional text transformation.

These Arabic-specific models, while valuable for their target tasks, operate
at scales (hundreds of millions of parameters) far below current instruction-tuned
LLMs.
More recent results have shown that large multilingual instruction-tuned models
such as GPT-4, Llama-3, and Qwen3 outperform Arabic-specific models on most
Arabic NLP benchmarks, primarily because the scale advantage outweighs the
benefit of Arabic-only pre-training.
Qwen3-4B was chosen for this study as the most capable openly accessible model
at a size deployable on free cloud GPU tiers.

% ============================================================
\section{Prompt Engineering}
\label{sec:prompting}
% ============================================================

\subsection{Zero-Shot Prompting}
\label{subsec:zero_shot}

In zero-shot prompting, the model receives an instruction and an input without
any demonstration examples \cite{brown2020language,liu2023pretrain}.
Performance depends entirely on the clarity of the instruction and the model's
pre-training coverage of the task.
For post-OCR correction, the zero-shot instruction must communicate both the
task (correct errors in the text) and the constraint (do not change what is
correct).

A key design decision is the degree of \emph{aggressiveness} in the instruction.
An aggressive prompt (``correct all errors in this text'') primes the model to
intervene on every uncertainty, producing rewrites of already-correct tokens.
A conservative prompt (``return the text unchanged if it is correct; fix only
clear errors'') primes the model to leave the text alone unless it is confident
of an error.
For a task where the vast majority of tokens in the input are already correct
(OCR baseline CER of 1.37\% on normal samples), conservative prompting is
the rational choice.
The design and evaluation of this choice is a central finding of Chapter~5.

\subsection{Few-Shot Prompting}
\label{subsec:few_shot}

Few-shot prompting, introduced in GPT-3 \cite{brown2020language}, provides
the model with a small number of input-output demonstration pairs before the
actual query.
This enables \emph{in-context learning}: the model infers the task format and
the expected transformation from the examples without any weight update.

Several practical decisions govern few-shot performance:

\begin{itemize}
  \item \textbf{Number of examples.}
        More examples generally improve performance up to the point where they
        saturate the context window.
        This study uses 5 examples, balancing context length against
        representativeness.

  \item \textbf{Example selection.}
        Random selection may draw all examples from the same error type.
        Diversity-maximising selection ensures that the examples span multiple
        error categories, giving the model a broader template for what
        corrections look like.

  \item \textbf{Source of examples.}
        The choice of few-shot corpus matters.
        QALB annotations cover human typing errors (omitted diacritics, wrong
        hamza form, case-agreement errors) that have a partial but imperfect
        overlap with OCR errors (visual confusion between similar-shape characters).
        This \emph{transfer gap} --- the mismatch between human-error and OCR-error
        distributions --- is a known limitation of few-shot approaches using
        human-authored error corpora for OCR correction.
\end{itemize}

\subsection{Chain-of-Thought and Instruction Augmentation}
\label{subsec:cot}

Chain-of-thought (CoT) prompting \cite{wei2022chain} instructs the model to
generate intermediate reasoning steps before producing the final answer.
CoT has been shown to improve performance on multi-step reasoning tasks such as
arithmetic, symbolic manipulation, and commonsense inference.

Chain-of-thought is not used in this study.
Post-OCR correction is a one-pass text-to-text transformation that does not
require multi-step reasoning: for each position in the input, the model either
copies the character or replaces it.
There is no decomposition into sub-problems that would benefit from intermediate
reasoning.
Furthermore, enabling chain-of-thought in Qwen3 (via the thinking mode) adds
substantial latency, unpredictable output length, and the risk of the reasoning
trace contaminating the extracted correction.
The conservative instruction in the system prompt already encodes the key
cognitive constraint (``prefer no change unless certain'') without requiring
explicit reasoning steps.

\subsection{Retrieval-Augmented Generation}
\label{subsec:rag}

Retrieval-Augmented Generation (RAG) \cite{lewis2020retrieval} combines a
parametric language model with a non-parametric retrieval component.
At inference time, a retrieval system queries a large external corpus with an
embedding of the input, retrieves the most similar passages, and provides these
as additional context in the prompt.
This allows the model to ground its outputs in specific textual evidence rather
than relying solely on knowledge encoded in its weights.

The RAG pipeline consists of the following components:

\begin{enumerate}
  \item \textbf{Corpus indexing.}
        The external corpus is split into sentences or passages, each embedded
        by a dense encoder into a fixed-dimensional vector.
        Embeddings are stored in a vector index that supports efficient
        nearest-neighbour search.

  \item \textbf{Query encoding.}
        The input text is embedded using the same encoder.

  \item \textbf{Retrieval.}
        The index is searched for the $k$ corpus items with embeddings most
        similar to the query, typically by inner product or cosine similarity.
        This study uses the FAISS \cite{johnson2019faiss} library with a flat
        inner-product index.

  \item \textbf{Prompt augmentation.}
        The retrieved passages are formatted and prepended to the system prompt
        as reference context.
\end{enumerate}

\noindent
This study uses the \texttt{paraphrase-multilingual-MiniLM-L12-v2} model
\cite{reimers2019sentence} for both corpus and query encoding.
This is a compact 12-layer multilingual sentence encoder with strong Arabic
support, suitable for embedding the 200,000-sentence OpenITI index on a CPU
without excessive memory usage.

An important consideration for RAG quality is \emph{domain match}: if the
retrieval corpus covers a different domain or time period than the input text,
retrieved passages may not be relevant and could mislead the model.
The domain mismatch between the classical OpenITI corpus and modern
newspaper text is a known limitation that was identified in preliminary
experiments and led to the exclusion of RAG from the final experimental
design of this thesis.

\subsection{Self-Reflective and Meta-Cognitive Prompting}
\label{subsec:self_reflective}

Self-reflective prompting refers to techniques that incorporate information
about the model's own behaviour into the prompt to guide future outputs.
The Reflexion framework \cite{shinn2023reflexion} demonstrates that language
agents can improve over time by receiving verbal feedback about their past
failures and incorporating this feedback into their decision-making.
Similarly, self-critique prompting asks the model to evaluate its own output
and revise it, akin to verbal reinforcement learning.

Phase~4 operationalises a lightweight variant of this idea.
Rather than asking the model to introspect on its current output at inference
time -- which adds latency and the risk of circular reasoning -- the model is
given \emph{pre-computed statistics} about its systematic failure patterns,
derived from analysing its corrections on a separate training split.
These statistics are converted to Arabic diagnostic sentences injected into
the system prompt.
The approach is a form of offline meta-learning: the model's systematic biases
are characterised before deployment, and this characterisation is used to
nudge its inference behaviour.

This design differs from standard Reflexion in two ways:
it does not require an execution environment or reward signal;
and it provides population-level statistics (``you fix insertion errors 97\% of
the time but introduce errors on 25\% of correct tokens'') rather than
sample-level feedback.
Whether this coarser form of self-knowledge translates into improved behaviour
at inference time is one of the empirical questions studied in this thesis.

% ============================================================
\section{Knowledge Sources Used in This Study}
\label{sec:knowledge_sources}
% ============================================================

\subsection{Confusion Matrix as a Knowledge Source}
\label{subsec:ks_confusion}

An OCR confusion matrix records, for each ground-truth character, the frequency
with which the OCR engine substituted other characters in its place.
It is a compact, data-driven summary of the engine's systematic recognition
errors.
Unlike generic Arabic spelling rules, a confusion matrix is engine-specific:
it reflects the particular visual features that this engine's neural network has
difficulty discriminating.

In this study, the confusion matrix is generated from Phase~1 analysis
(Section~\ref{sec:datasets} and Chapter~4) and is used in Phase~3 to inject
the top-10 most frequent confusion pairs into the system prompt.
The hypothesis is that naming these specific pairs directs the model's attention
to the character positions most likely to contain errors.

\subsection{Arabic Orthographic Rules}
\label{subsec:ks_rules}

Arabic orthographic rules govern the correct spelling of characters that OCR
commonly misidentifies.
Relevant rule categories include: hamza placement (al-qat' vs.\ al-wasl),
taa marbuta vs.\ ha, alef maqsura vs.\ ya, tanwin placement, and definite
article assimilation.
These rules are grounded in classical Arabic grammar and are taught in Arabic
primary education.

Unlike the confusion matrix (which is descriptive --- ``this engine makes this
mistake''), orthographic rules are prescriptive --- ``the correct form is X''.
Both types of knowledge can in principle direct the model's attention to
correction-relevant phenomena; however, preliminary experiments in this study
found that rule-augmented prompting did not improve on zero-shot correction
(the rules tended to duplicate guidance already embedded in the base prompt),
and this approach was excluded from the final experimental design.
The confusion-matrix injection in Phase~3 proved the more effective descriptive
alternative.

\subsection{The QALB Corpus}
\label{subsec:ks_qalb}

The Qatar Arabic Language Bank (QALB) was developed for the QALB Shared Task
on Automatic Arabic Text Correction \cite{mohit2014first,zaghouani2014large}.
It contains thousands of human-authored Arabic sentences from online news
comments, annotated with corrections for spelling, grammatical, punctuation,
and other errors by trained annotators.

QALB has been used in related work as a source of few-shot demonstration pairs
for Arabic correction tasks.
After filtering for OCR-relevance and length, several hundred pairs are available
for potential use.
The primary limitation of QALB as an OCR correction resource is the
\emph{transfer gap}: QALB errors are human typing errors, while OCR errors are
visual confusion errors.
Human typers tend to make phonological errors (writing a word as it sounds in
dialectal pronunciation) or agreement errors (wrong gender or case marking),
while OCR engines make visual confusion errors (confusing characters with similar
shapes).
The overlap between these two error distributions is partial, and the few-shot
examples do not perfectly model what the model is being asked to correct.

\subsection{The OpenITI Corpus}
\label{subsec:ks_openiti}

The Open Islamicate Texts Initiative (OpenITI) is a large-scale digitised
collection of texts from the Islamicate scholarly tradition, containing classical
Arabic, Persian, and other languages \cite{nigst2020openiti,romanov2017openiti}.
With hundreds of millions of words of Arabic text spanning many centuries and
genres, it provides a large pool of grammatically correct Arabic sentences for
dense retrieval.

OpenITI represents a potential retrieval corpus for RAG-based correction
approaches. Its primary limitation is \emph{domain mismatch}: OpenITI
texts are classical and historical, while PATS-A01 documents are modern
newspaper text. The semantic and lexical overlap between medieval religious
scholarship and contemporary news reporting is low. Preliminary experiments
with OpenITI-based RAG showed severe degradation due to this mismatch, and
the RAG approach was not included in the final experimental design of this
thesis. A retrieval corpus drawn from modern Arabic newspapers would be
required for effective RAG-based correction.

\subsection{CAMeL Tools}
\label{subsec:ks_camel}

CAMeL Tools \cite{obeid2020camel} is an open-source Python toolkit for Arabic
NLP, developed at the Computational Approaches to Modeling Language Lab at
NYU Abu Dhabi.
The components used in this study are:

\textbf{MorphologicalAnalyzer.}
Given an Arabic word token, the analyser queries the MSA morphological database
(\texttt{calima-msa-r13}) and returns all valid morphological analyses:
each analysis specifies the lemma, root, word pattern, part-of-speech, and
inflectional features (gender, number, case, tense, etc.).
A word with zero analyses is morphologically invalid -- it is not a recognisable
MSA word form.

\textbf{Application in Phase~1.}
CAMeL Tools is used in Phase~1 to classify OCR error tokens as either
\emph{non-words} (no valid morphological analysis) or \emph{valid-but-wrong}
(at least one valid analysis but not the correct one in context).
The distinction is informative: if most OCR errors are valid-but-wrong, then
dictionary lookup cannot detect them and context is required for correction.

\textbf{Application in Phase~5.}
The \emph{revert strategy} described in Chapter~3 uses the morphological analyser
to post-process LLM corrections.
Words introduced by the LLM that fail morphological analysis while the original
OCR token passed are reverted to the OCR original.

% ============================================================
\section{Related Work on Post-OCR Correction}
\label{sec:related}
% ============================================================

\subsection{Classical Approaches}
\label{subsec:related_classical}

The earliest computational approaches to post-OCR correction relied on lexicons
and language models to identify and resolve implausible recognitions.
Tong and Evans \cite{tong1996statistical} proposed a statistical approach using
a unigram word model and a character-level confusion model derived from known
OCR error patterns.
Their method identified likely errors by flagging tokens not found in the lexicon
or whose recognition score was below a threshold, then proposed corrections by
searching the lexicon for candidates within a small edit distance.
This work established the noisy channel framework for OCR correction that
remained influential for over a decade.

Kolak and Resnik \cite{kolak2005ocr} developed this framework further with an
explicit noisy channel model combining character-level OCR confusion probabilities
with a word-level language model.
Their system formulated post-OCR correction as a decoding problem: given the
OCR output, find the most probable original text under the joint model.
This framing is conceptually equivalent to the machine translation noisy channel
model and shares its strengths (principled probabilistic framework) and
limitations (requires domain-specific language models and large training corpora
of OCR errors).

The limitation common to all classical approaches is their reliance on
fixed-vocabulary lexicons and short-context language models.
They cannot resolve \emph{valid-but-wrong} errors -- tokens that are real words
but not the intended word in context -- without long-range context, and they
require error-specific training data for each OCR engine being corrected.

\subsection{Neural Sequence-to-Sequence Correction}
\label{subsec:related_neural}

The deep learning era brought encoder-decoder models to post-OCR correction.
Dong and Smith \cite{dong2018multi} applied a multi-input attention model that
combines character-level and word-level representations of the OCR output,
enabling the model to attend to both local character confusion patterns and
broader word-level context.
Their system significantly outperformed classical methods on historical English
document collections, demonstrating that neural sequence models can implicitly
learn OCR error distributions from paired training data.

Transformer-based sequence-to-sequence models have since been applied to
post-OCR correction with further improvements.
The ICDAR 2019 Post-OCR Correction Competition \cite{rigaud2019icdar} attracted
systems based on BERT, character-level transformers, and ensemble methods.
These systems substantially improved over the OCR baseline on historical French
and English text.

The fundamental challenge for these approaches in the Arabic context is the
\emph{scarcity of training data}.
Neural sequence-to-sequence correction requires large quantities of paired
(OCR output, ground truth) examples to learn the error model for a specific
engine.
Such data is relatively abundant for historical English and French (where
digitised collections with corresponding scans and transcripts exist), but
scarce for modern Arabic.
Creating a training set for Arabic OCR correction would require either manual
transcription of thousands of document scans or a data augmentation strategy
that synthetically generates OCR errors.
This data requirement motivates the zero-shot and prompt-based approaches
investigated in this thesis, which require no paired training data.

\subsection{LLM-Based Post-OCR Correction}
\label{subsec:related_llm}

The emergence of large instruction-tuned language models created a new paradigm
for post-OCR correction: the model is instructed to correct errors in the OCR
output without any task-specific fine-tuning.
Kang et al.\ \cite{kang2024llmocr} demonstrated that ChatGPT-level models are
effective at English OCR post-correction in a zero-shot setting, achieving
substantial CER reductions on several historical English document benchmarks.
Their results show that the implicit knowledge about spelling, grammar, and
vocabulary encoded in pre-training is sufficient to correct many OCR errors
without any paired training data.

This line of work directly motivates the present study.
However, Kang et al.\ and similar work focus entirely on English.
Whether the same approach generalises to Arabic -- a morphologically richer
script-level script with a different error distribution -- is not addressed.
Furthermore, no prior work has compared multiple prompt augmentation strategies
under controlled conditions for Arabic OCR correction.

The phenomenon of \emph{over-correction} -- where the LLM changes tokens that
the OCR rendered correctly, introducing errors rather than fixing them -- is
discussed qualitatively in some LLM correction work but has not been
systematically studied as the dominant failure mode in the Arabic setting.

\subsection{Arabic-Specific OCR Correction}
\label{subsec:related_arabic}

Work specifically targeting Arabic OCR correction is limited.
Several groups have applied AraBERT fine-tuned on QALB-like data to Arabic
error correction, achieving competitive results on the QALB shared task.
However, the QALB task targets human writing errors in online text comments
rather than OCR-specific visual confusion errors, and fine-tuned models trained
on QALB data are not well-calibrated for the OCR error distribution.

Rule-based approaches using Arabic spell-checkers such as Hunspell with Arabic
dictionaries have also been applied to post-processing Arabic OCR output.
These systems can flag non-word tokens (similar to classical approaches) but
cannot resolve valid-but-wrong errors.

To the best of the author's knowledge, no published study has:
(a) systematically compared multiple distinct prompt augmentation strategies
for Arabic post-OCR correction in an isolated experimental design;
(b) applied self-reflective prompting to the OCR correction task in any language;
or (c) characterised the over-correction failure mode in the Arabic post-OCR
setting with quantitative evidence.

\subsection{Positioning of This Work}
\label{subsec:positioning}

This thesis makes the following contributions relative to the existing literature:

\begin{enumerate}
  \item \textbf{Systematic multi-dataset evaluation of augmentation strategies.}
        Prior work on LLM-based OCR correction has at most compared zero-shot
        with one or two augmentation conditions on a single dataset.
        This study isolates three isolated knowledge augmentation conditions
        (confusion-matrix injection, self-reflective prompting, and morphological
        post-processing), tests all pairwise combinations with ablation, and
        evaluates across nine datasets (eight typewritten fonts and handwritten
        text) using the same model and evaluation protocol.

  \item \textbf{Self-reflective prompting for OCR correction.}
        The Phase~4 strategy --- injecting pre-computed statistics about the
        model's own training-split error patterns into the inference prompt ---
        is, to the author's knowledge, the first application of self-reflective
        prompting to OCR correction in any language.

  \item \textbf{Quantitative study of over-correction in Arabic OCR.}
        The finding that LLM interventions harm fonts with CER below approximately
        8\%, and that the over-correction rate grows as more knowledge is injected
        into the prompt, is a novel empirical characterisation of a failure mode
        that deserves attention from practitioners deploying LLMs for text
        normalisation tasks.

  \item \textbf{Arabic-specific multi-font evaluation.}
        The PATS-A01 corpus provides eight typewritten Arabic fonts with an
        aligned train/validation split, enabling precise attribution of CER
        changes to specific prompt interventions while controlling for content
        effects across fonts.
\end{enumerate}
